% =========================================================================
%  Estilos y macros gráficas comunes del Tema 2
%  - \mapcolor{t}: escala cromática «campo» (t en [0,1]); define el color mapc
%  - Estilos TikZ heredados de los temas 0 y 1 (ejes, vectores, versores...)
%  - \panelMMA: marco de las figuras clonadas de Mathematica (StreamPlot,
%    ContourPlot y VectorPlot) que dibuja los datos de figuras/datos/*.tex
% =========================================================================

% --- Escala cromática (5 tramos entre cm0 y cm4) ------------------------
\newcommand{\mapcolor}[1]{%
  \pgfmathsetmacro{\mt}{min(max(#1,0),0.9999)*4}%
  \pgfmathtruncatemacro{\mi}{floor(\mt)}%
  \pgfmathtruncatemacro{\mj}{\mi+1}%
  \pgfmathsetmacro{\mf}{100*(\mt-\mi)}%
  \colorlet{mapc}{cm\mj!\mf!cm\mi}}

% --- Estilos TikZ generales ------------------------------------------------
\tikzset{
  >=Stealth,
  eje/.style={->, line width=0.6pt, grispizarra},
  ejeprof/.style={->, line width=0.7pt, magprof},
  vect/.style={->, line width=1pt, magprof},
  vectrojo/.style={->, line width=1pt, rojoteo},
  vectverde/.style={->, line width=1pt, verdeteo},
  fino/.style={line width=0.4pt, grispizarra},
  punteado/.style={densely dotted, line width=0.5pt, grispizarra},
  etiq/.style={font=\small, text=black!85},
  cuadro/.style={draw=grispizarra!60, line width=0.4pt, fill=white, inner sep=2pt},
  trayecto/.style={line width=1.1pt, rojoteo},
  orienta/.style={decoration={markings, mark=at position 0.55 with {\arrow{Stealth[length=5pt]}}}, postaction={decorate}},
  % líneas de campo y equipotenciales en la paleta del documento
  lcampo/.style={magacero, line width=0.55pt, postaction={decorate},
                 decoration={markings, mark=at position 0.5 with {\arrow{Stealth[length=3.6pt]}}}},
  lcampoinv/.style={magacero, line width=0.55pt, postaction={decorate},
                 decoration={markings, mark=at position 0.5 with {\arrowreversed{Stealth[length=3.6pt]}}}},
  lequi/.style={pizarra!70, line width=0.45pt, densely dashed},
}

% --- Estilo base de ejes pgfplots (como en los temas anteriores) -----------
\pgfplotsset{
  tema/.style={
    tick label style={font=\scriptsize}, label style={font=\small},
    axis lines=left, axis line style={grispizarra},
    grid=major, grid style={grisclaro, line width=0.4pt},
    every tick/.style={grispizarra},
    legend style={font=\scriptsize, draw=none, fill=white, fill opacity=0.85,
                  text opacity=1},
    /pgf/number format/use comma,
  },
  % Ejes con marco completo, como los «Plot» de Mathematica de las figuras del tema
  marcoMMA/.style={
    tick label style={font=\scriptsize}, label style={font=\small},
    axis line style={grispizarra!80, line width=0.5pt},
    every tick/.style={grispizarra}, tick align=inside,
    grid=major, grid style={grisclaro!70, line width=0.3pt},
    /pgf/number format/use comma,
    legend style={font=\scriptsize, draw=none, fill=white, fill opacity=0.85,
                  text opacity=1},
  },
  mapa/.style={
    tick label style={font=\scriptsize}, label style={font=\small},
    axis line style={grispizarra}, every tick/.style={grispizarra},
    /pgf/number format/use comma,
    colormap name=campo,
  },
}

% =========================================================================
%  Paneles «tipo Mathematica» (ventana [-L,L]^2 en unidades de \ell)
%  Uso dentro de un tikzpicture con x=y=<escala>:
%    \panelMMA{L}{fondo}{radio del disco central}{contenido}
%  El contenido se recorta a la ventana; luego se dibujan el disco blanco
%  (región excluida, donde están las fuentes) y el marco con marcas.
% =========================================================================
% Estilos de los objetos que escribe scripts/generar_datos.py
\tikzset{
  mmaLC/.style={mmalinea, line width=0.42pt},
  mmaEQ/.style={mmaequi, line width=0.28pt},
  mmaEQsolo/.style={black, line width=0.4pt},
  mmaFL/.style={-{Stealth[length=3.3pt, width=2.7pt]}, mmalinea, line width=0.42pt},
}
% Macros por defecto (cada figura puede redefinirlas antes de \input)
\newcommand{\EQ}[1]{\draw[mmaEQ] plot coordinates {#1};}
\newcommand{\LC}[1]{\draw[mmaLC] plot coordinates {#1};}
\newcommand{\FL}[3]{\draw[mmaFL] (#1,#2) -- ++(#3:0.02);}
\newcommand{\VF}[5]{%
  \mapcolor{#5}%
  \pgfmathsetmacro{\hl}{2.0+2.4*#5}\pgfmathsetmacro{\lw}{0.32+0.45*#5}%
  \draw[-{Stealth[length=\hl pt,width=0.85*\hl pt]}, line width=\lw pt, mapc]
     ({#1-0.5*#3},{#2-0.5*#4}) -- ++(#3,#4);}

% Marco con marcas en -L, -L/2, 0, L/2, L y rótulos de los ejes
% \marcoMMA{L}{paso de marcas}{rótulo x}{rótulo y}
\colorlet{marcoc}{mmamarco}
\newcommand{\marcoMMA}[4]{%
  \draw[marcoc, line width=0.55pt] (-#1,-#1) rectangle (#1,#1);
  \foreach \t in {-#1,-#2,0,#2,#1}{%
    \draw[grispizarra, line width=0.35pt] (\t,-#1) -- ++(0,0.025*#1*2);
    \draw[grispizarra, line width=0.35pt] (-#1,\t) -- ++(0.025*#1*2,0);
    \draw[grispizarra, line width=0.35pt] (\t,#1) -- ++(0,-0.025*#1*2);
    \draw[grispizarra, line width=0.35pt] (#1,\t) -- ++(-0.025*#1*2,0);
    \node[below=1pt, font=\tiny, inner sep=1pt] at (\t,-#1) {$\pgfmathprintnumber{\t}$};
    \node[left=1pt, font=\tiny, inner sep=1pt] at (-#1,\t) {$\pgfmathprintnumber{\t}$};}
  \node[font=\small, anchor=north] at (0,-#1) [yshift=-9pt] {#3};
  \node[font=\small, anchor=south, rotate=90] at (-#1,0) [xshift=0pt, yshift=12pt] {#4};}

% \panelMMA{L}{color de fondo}{r0}{contenido}
\newcommand{\panelMMA}[4]{%
  \fill[#2] (-#1,-#1) rectangle (#1,#1);
  \begin{scope}\clip (-#1,-#1) rectangle (#1,#1); #4 \end{scope}
  \ifdim #3pt>0pt
    \fill[white] (0,0) circle (#3);
    \draw[mmamarco!80, line width=0.4pt] (0,0) circle (#3);
  \fi}

% Vector dipolar rojo en el centro de un panel (altura ±h)
\newcommand{\dipoloRojo}[1]{%
  \draw[-{Stealth[length=4pt,width=3.6pt]}, mmarojo, line width=1.3pt] (0,-#1) -- (0,#1);}

% =========================================================================
%  Primitivas 3D para tdplot (vista fija: tdplotsetmaincoords{70}{125}),
%  heredadas del Tema 1.
% =========================================================================
\def\VA{125}
% \tdCirc[opciones]{radio}{z}   (mitad frontal continua, trasera a trazos)
\newcommand{\tdCirc}[3][]{%
  \draw[#1, domain=\VA-180:\VA, samples=36, variable=\tdt]
    plot ({#2*cos(\tdt)},{#2*sin(\tdt)},{#3});
  \draw[#1, densely dashed, domain=\VA:\VA+180, samples=36, variable=\tdt]
    plot ({#2*cos(\tdt)},{#2*sin(\tdt)},{#3});}
% \tdEsfera{radio}{color}  (esfera sombreada con ecuador)
\newcommand{\tdEsfera}[2]{%
  \begin{scope}[tdplot_screen_coords]
    \shade[ball color=rosaclaro!80, opacity=0.45] (0,0) circle (#1);
    \draw[#2, line width=0.7pt] (0,0) circle (#1);
  \end{scope}
  \tdCirc[#2!60, line width=0.4pt]{#1}{0}}

% =========================================================================
%  Región de carga V' para los diagramas geométricos (secciones 3 y 4):
%  contorno suave, separado de los ejes, con sombreado radial tenue.
%  \regionV{desplazamiento}  (el contorno está centrado cerca de (2.25,2.7))
% =========================================================================
\def\contornoV{plot[smooth cycle, tension=0.75] coordinates
  {(0.95,2.5) (1.35,3.5) (2.45,3.95) (3.55,3.45) (3.6,2.25) (2.6,1.55) (1.45,1.65)}}
\newcommand{\regionV}{%
  \shade[inner color=rosaclaro!25, outer color=rosaclaro!90] \contornoV;
  \draw[magacero, line width=0.75pt] \contornoV;}
% Ejes cartesianos en perspectiva caballera con el origen marcado
% \ejesO{long z}{long y}{long x}
\newcommand{\ejesO}[3]{%
  \draw[eje] (0,0) -- (0,#1) node[above, etiq] {$z$};
  \draw[eje] (0,0) -- (#2,0) node[right, etiq] {$y$};
  \draw[eje] (0,0) -- ({-0.76*#3},{-0.65*#3}) node[below left=-2pt, etiq] {$x$};}
% Elemento de volumen (cubito en perspectiva) centrado en #1
\newcommand{\cubito}[1]{%
  \draw[magprof, line width=0.5pt, fill=white] ($(#1)+(-0.12,-0.12)$) rectangle ($(#1)+(0.12,0.12)$);
  \draw[magprof, line width=0.5pt, fill=rosaclaro!40] ($(#1)+(-0.12,0.12)$) -- ++(0.07,0.07) -- ++(0.24,0) -- ++(-0.07,-0.07) -- cycle;
  \draw[magprof, line width=0.5pt, fill=rosaclaro!70] ($(#1)+(0.12,0.12)$) -- ++(0.07,0.07) -- ++(0,-0.24) -- ++(-0.07,-0.07) -- cycle;}
